\subsection{Background and Motivation}

Finding a physical law from measurements is one of the oldest tasks in science. A researcher collects data, guesses a relationship, tests it, and keeps improving the guess until it holds. For a long time this could only be done by people. That started to change when computers were first used to search for equations directly. Schmidt and Lipson (2009) showed that an algorithm could recover laws such as energy conservation from motion data without being told any physics beforehand.

Since then, artificial intelligence has taken on larger parts of scientific work. AlphaFold predicted protein structures at an accuracy close to experiments (Jumper et al., 2021), and a review by Wang et al. (2023) describes how AI now helps scientists form hypotheses, design experiments and interpret large datasets. More recently, large language models (LLMs) have been used to carry out the whole research cycle. The AI Scientist, for example, generates ideas, runs experiments and writes up the results by itself (Lu et al., 2024).

These systems raise a question that is hard to answer with ordinary tests. When an LLM gives the correct formula for a physical law, did it work the law out from the data, or did it remember the formula from its training? Well known laws like Hooke's law or Coulomb's law appear in countless textbooks, so a correct answer does not prove that any discovery took place. To test discovery properly, the law has to be one the model has not seen, and the model should have to collect the evidence itself as a scientist would.

Our project is motivated by this problem. We redesigned familiar physics laws so that their textbook forms are wrong, placed them inside a simulated laboratory, and asked LLM agents to find them by running their own experiments. There is also a practical side. Large models are costly to run, so it matters whether a small model can reach the same level by improving its own strategy over time. If it can, automated discovery becomes much cheaper. If it cannot, that tells us where the real limits of these models are.
